The quotients of Theorems~\ref{thm:upper-unit-quotient} and~\ref{thm:upper-idempotent-quotient} coincide and have a particularly
explicit form in the following broad class. Call an idempotent support semiring \emph{integral} here if its
multiplicative unit is its greatest element. This use of the word does
not mean absence of zero divisors.

\begin{theorem}[removing transition kernels]
\label{thm:integral-support-reduction}
Let $T=(Q,F)$ be a unital inverse semiring with integral support $Q$.
Put $R=F(1)$ and $I_q=\im(F(q)\to R)$. Then the $I_q$ form an
increasing diagram of two-sided ideals of the unital ring $R$, with
$I_1=R$ and $I_qI_p\subseteq I_{qp}$. Their inclusions and ordinary
ring products define an inverse semiring $T_{\mathrm{sep}}=(Q,I)$ and
a surjective homomorphism $T\to T_{\mathrm{sep}}$ that preserves every
support. This is precisely either of these two zero-detection quotients.

For every idempotent $n\times n$ support matrix $E$, $n\geq1$, the
additive matrix fiber $A_E=\prod_{i,j}F(E_{ij})$ is a possibly
nonunital ring. The kernel $N_E$ of its map to
$\prod_{i,j}I_{E_{ij}}$ satisfies
\[
 N_E^{\,n+1}=0.
\]
Consequently extension of scalars induces a bijection on isomorphism
classes of finitely generated projective left semimodules, preserving
direct sums. In particular,
\[
 V(T)\simeq V(T_{\mathrm{sep}}),\qquad
 K_0(T)\simeq K_0(T_{\mathrm{sep}}).
\]
There is no finiteness, height or commutativity hypothesis on $Q$, and
no finiteness or commutativity hypothesis on the coefficient groups or
the unit ring. The assertion concerns projective isomorphism classes;
it does not assert equivalence of the full projective categories.
\end{theorem}

For integral support, transition kernels therefore create no additional
finite-projective classes. The residual coefficient data are actual
ideals in one ordinary ring. Projectives over that ideal diagram still
require classification; the preceding theorem retains its unrestricted scope.
